\Askhsh{23215}{4}
\begin{ylh}{1}
    \item 1.2 Συναρτήσεις
    \item 1.3 Μονότονες συναρτήσεις - Αντίστροφη συνάρτηση
    \item 1.8 Συνέχεια συνάρτησης
    \item 2.5 Θεώρημα Μέσης Τιμής Διαφορικού Λογισμού
    \item 2.6 Συνέπειες του Θεωρήματος Μέσης Τιμής
    \item 2.7 Τοπικά ακρότατα συνάρτησης.
\end{ylh}
Δίνεται συνάρτηση $f : \mathbb{R} \to \mathbb{R}$ για την οποία ισχύει $f'(x) \neq 0$ για κάθε $x \in \mathbb{R}$.

\begin{alist}
\item Να αποδείξετε ότι η συνάρτηση $f$ είναι 1-1. \monades{5}

Δίνεται επιπλέον ότι
\begin{itemize}
    \item η συνάρτηση $f'$ είναι συνεχής,
    \item $f(0) = -1$ και $f(2) = 1$.
\end{itemize}

\item Να αποδείξετε ότι υπάρχει εφαπτομένη της γραφικής παράστασης της $f$ που είναι παράλληλη στην ευθεία $y=x$. \monades{5}

\item
\begin{rlist}
    \item Να αποδείξετε ότι η συνάρτηση $f$ είναι γνησίως αύξουσα. \monades{4}
    \item Να αποδείξετε ότι η γραφική παράσταση της $f$ τέμνει τον άξονα $x'x$ σε ένα μόνο σημείο με τετμημένη $x_0 \in (0,2)$. \monades{5}
\end{rlist}

\item Αν $g$ είναι μια συνάρτηση για την οποία ισχύει ότι $g'(x) = -f(x)$ για κάθε $x \in \mathbb{R}$ να αποδείξετε ότι η $g$ παρουσιάζει ολικό μέγιστο στο $x_0$. \monades{6}
\end{alist}